\begin{helper}
Let \(V\) be finite and let \(a,b,c\in V\) be pairwise distinct.  Let
\(E\subseteq\mathbb R^V\) be Lebesgue measurable for the finite product
Lebesgue measure, and let \(g:\mathbb R^3\to\mathbb R\) be continuous.  Assume
that \(g(x,y,z)\ge0\) whenever \(0\le x\le y\le z\le1\), and assume that every
\(q\in E\) satisfies
\[
0\le q_a\le q_b\le q_c\le 1.
\]
Assume also that for every \(0\le x\le y\le z\le1\), the slice of \(E\) with
\(q_a=x\), \(q_b=y\), and \(q_c=z\), viewed as a subset of the remaining
coordinates \(V\setminus\{a,b,c\}\), has product Lebesgue outer measure at
most \(\operatorname{ofReal}(g(x,y,z))\).  Then
\[
\operatorname{vol}(E)\le
\operatorname{ofReal}\left(\int_0^1\int_0^z\int_0^y
g(x,y,z)\,dx\,dy\,dz\right).
\]
\end{helper}
\begin{proof}
Because \(a,b,c\) are pairwise distinct, the coordinate map
\(\mathbb R^V\to\mathbb R^3\times \mathbb R^{V\setminus\{a,b,c\}}\) sending
\(q\) to \(((q_a,q_b,q_c),(q_v)_{v\notin\{a,b,c\}})\) is a permutation of the
finite product coordinates and preserves product Lebesgue measure.  Apply
Tonelli--Fubini to the measurable set \(E\) under this coordinate
identification.  The containment hypothesis says that the indicator of \(E\)
vanishes unless the distinguished coordinates lie in the simplex
\(0\le x\le y\le z\le1\).  Thus \(\operatorname{vol}(E)\) is bounded by the
iterated integral, over this simplex, of the outer product measure of the
corresponding outside-coordinate slice.

For each point of the simplex the slice hypothesis bounds that outside measure
by \(\operatorname{ofReal}(g(x,y,z))\).  The continuity of \(g\) supplies the
measurability and interval-integrability needed for the displayed iterated
Lebesgue integrals, and the assumed nonnegativity on the simplex justifies
passing between the nonnegative outer-measure bound and
\(\operatorname{ofReal}(g(x,y,z))\).  Monotonicity of the nonnegative integral
therefore gives the displayed upper bound.
\end{proof}
